\chapter{Partitioned Point-to-Point Communication}
\mpitermtitleindexsubmain{partitioned point-to-point}{communication}
\label{sec:part}
\label{chap:part}

\section{Introduction}
\label{sec:part-intro}
Partitioned communication extends persistent point-to-point communication as  defined in
Chapter~\ref{chap:pt2pt}. Partitioned communication operations are matched based on the order in which the local initialization calls are performed. Partitioned
communication is ``partitioned'' because it allows for multiple contributions of
data to be made, potentially, from multiple actors (e.g., threads or tasks) in an \MPI/ process to a single communication
operation.

\begin{users}
The techniques of partitioned communication were known as ``finepoints'' 
before their adoption into the \MPI/ standard. 
We refer the interested reader to the original literature describing the
design goals, functioning, initial implementation 
and performance improvements~\cite{finepoints_isc,finepoints}.
\end{users}

Partitioned communication operations use
a persistent communication style that involves a sequence of start and test or wait operations. For this
sequence, partitioned communications use \mpifunc{MPI\_START} or \mpifunc{MPI\_STARTALL} calls
and completion mechanisms (e.g., \mpifunc{MPI\_TEST} or \mpifunc{MPI\_WAIT}).
Partitioned communication is different in three fundamental ways from persistent
point-to-point operations in \MPI/. 
First, partitioned communication
allows additional partitioned test function calls that can expose partial completion
of the operation.
Second, partitioned communication may perform all
of the initialization required to enable data transfer as
early as its initialization phase. Third, partitioned communication allows for
\MPI/ to be independently notified of multiple contributions from the 
send-side to a single data buffer of a single \MPI/ message.

\begin{rationale}
The rationale behind having different initialization behavior allowed for
partitioned communication as opposed to persistent point-to-point communication is to
enable flexibility and optimization possibilities in implementations.
Buffer setup can occur in the
partitioned communication initialization functions (see Section~\ref{sec:part_comm_init}). 
However, such negotiation can be deferred until
data is to be moved between two processes. This means that partitioned communication can
lazily negotiate as late as testing for completion of the operation on the first
iteration of a sequence of partitioned communication start and test or wait operations.
Matching still occurs as if matching happened at the partitioned communication initialization functions
as noted in the function descriptions.
\end{rationale}

\section{Semantics of Partitioned Point-to-Point Communication}
\mpitermtitleindex{semantics!partitioned point-to-point communication}
\label{sec:part-semantics} 
\MPI/ guarantees certain general properties of
partitioned point-to-point communication progress,
which are described in this section.

Persistent communications use opaque \mpiarg{MPI\_REQUEST} objects
as described in Section~\ref{sec:pt2pt}. Partitioned communication
uses these same semantics for \mpiarg{MPI\_REQUEST} objects.

Partitioned communication provides fine-grained transfers on either or
both sides of a send-receive operation described by requests.
Persistent communication semantics are ideal for partitioned
communication: they provide \mpifunc{MPI\_PSEND\_INIT} and
\mpifunc{MPI\_PRECV\_INIT} functions that allow partitioned
communication setup to occur prior to message transfers. Partitioned
communication initialization functions are local.  The partitioned
communication initialization includes inputs on the number of
user-visible partitions on the send-side and receive-side, which may
differ.  Valid partitioned communication operations must have
one or more partitions specified.

Once an \mpifunc{MPI\_PSEND\_INIT} call has been made, the user may start the
operation with a call to a starting procedure and complete the
operation with one \mpifunc{MPI\_PREADY} call for every
send partition followed by a call to a completing
procedure. A call to \mpifunc{MPI\_PREADY} notifies the
\MPI/ library that a specified portion of the data buffer (a specific partition) is ready to be sent.
Notification of partial completion can be done
via fine-grained \mpifunc{MPI\_PARRIVED} calls at the receiver before a final
\mpifunc{MPI\_TEST}/\mpifunc{MPI\_WAIT} on the request itself; the latter
represents overall operation completion upon success.  
A full set of methods
for starting and completing partitioned communication is given in the
following sections.

\begin{users}
Having a large number of receiver-side partitions can increase overheads 
as the completion mechanism may need to work with finer-grained notifications. Using a small
number of receiver-side partitions \emph{may} provide higher performance.

A large number of sender-side partitions may be aggregated by an \MPI/ implementation,
making performance concerns of a large number of sender-side partitions potentially less impactful
than receiver-side granularity. 
\end{users}

\begin{implementors}
It is expected that an \MPI/ implementation will attempt to balance
latency and aggregation for
data transfers for the requested partition counts on the sender-side and receiver-side
to allow optimization for different hardware. A high quality implementation may 
perform significant optimizations to enhance performance in this way; they may, for example,
resize the data transfers of the partitions to combine partitions in fractional
partition sizes (e.g., 2.5 partitions in a single data transfer). 
\end{implementors}

 
Example~\ref{part-example1} shows a simple partitioned transfer
in which the sender-side and receiver-side partitioning is identical in partition count.

\begin{example}
Simple partitioned communication example.
\label{part-example1}%
\exindex{Partitioned communication!Simple example}%
\Exindex{C}{Partitioned communication}{MPI\_Psend\_init,MPI\_Precv\_init,MPI\_Pready}%
%%HEADER
%%LANG: C
%%DECL: void error(const char *);
%%ENDHEADER
\begin{lstlisting}[language={[MPI]C}]
#include <stdlib.h>
#include "mpi.h"
#define PARTITIONS 8
#define COUNT 5
int main(int argc, char *argv[])
{
  double message[PARTITIONS*COUNT];
  int partitions = PARTITIONS;
  int source = 0, dest = 1, tag = 1, flag = 0;
  int myrank, i;
  MPI_Request request;
  MPI_Init(&argc, &argv);
  MPI_Comm_rank(MPI_COMM_WORLD, &myrank);
  if (myrank == 0)
  {
     MPI_Psend_init(message, partitions, COUNT, MPI_DOUBLE, dest, tag,
                    MPI_COMM_WORLD, MPI_INFO_NULL, &request);
     MPI_Start(&request);
     for(i = 0; i < partitions; ++i)
     {
        /* compute and fill partition #i, then mark ready: */
        MPI_Pready(i, request);
     }
     while(!flag)
     {
        /* do useful work #1 */
        MPI_Test(&request, &flag, MPI_STATUS_IGNORE);
        /* do useful work #2 */
     }
     MPI_Request_free(&request);
  }
  else if (myrank == 1)
  {
     MPI_Precv_init(message, partitions, COUNT, MPI_DOUBLE, source, tag,
                    MPI_COMM_WORLD, MPI_INFO_NULL, &request);
     MPI_Start(&request);
     while(!flag)
     {
        /* do useful work #1 */
        MPI_Test(&request, &flag, MPI_STATUS_IGNORE);
        /* do useful work #2 */
     }
     MPI_Request_free(&request);
  }
  MPI_Finalize();
  return 0;
}
\end{lstlisting}
\end{example}

\begin{rationale}
Partitioned communication is designed to provide opportunities for \MPI/ implementations to
optimize data transfers. \MPI/ is free to choose how many transfers to do within a partitioned communication send
independent of how many partitions are reported as ready to \MPI/ through
\mpifunc{MPI\_PREADY} calls.
Aggregation of partitions is permitted but not required.  Ordering of partitions  is permitted but not required.
A naive implementation can simply
wait for the entire message buffer to be marked ready before any transfer(s) occur and
could wait until the completion function is called on a request before transferring data.
However, this modality of communication gives \MPI/ implementations far more flexibility in data
movement than
nonpartitioned communications.
\end{rationale}

\subsection{Communication Initialization and Starting with Partitioning}
\label{sec:part_comm_init}
%\mpitermtitleindexmainsub{initialization}
%\mpitermtitleindexmainsub{starting}
%
Initialization of partitioned communication operations use the initialization calls described below.
Subsequent to initialization, \mpifunc{MPI\_START}/\mpifunc{MPI\_STARTALL} are used as the first indication
to \MPI/ that a message transfer will occur.
For send-side operations,
neither initializing nor starting the operation enables 
transfer of any part of the user buffer.
Freeing or canceling a partitioned communication request that is active (i.e., initialized and
started) and not completed is erroneous.
After the partitioned communication operation is started, individual partitions of a 
message are indicated as ready to be sent by \MPI/ via the \mpifunc{MPI\_PREADY}
function, described below.

\cdeclindex{MPI\_Request}
\begin{mpi-binding}
    function_name("MPI_Psend_init")
    parameter("buf", "BUFFER", direction="IN", desc=r"initial address of send buffer", constant="true")
    parameter("partitions", "PARTITION", direction="IN", desc=r"number of partitions")
    parameter("count", "XFER_NUM_ELEM_NNI", direction="IN", desc=r"number of elements sent per partition")
    parameter("datatype", "DATATYPE", direction="IN", desc=r"type of each element")
    parameter("dest", "RANK", direction="IN", desc=r"rank of destination")
    parameter("tag", "TAG", direction="IN", desc=r"message tag")
    parameter("comm", "COMMUNICATOR", direction="IN", desc=r"communicator")
    parameter("info", "INFO", desc=r"info argument")
    parameter("request", "REQUEST", direction="OUT", desc=r"communication request")
\end{mpi-binding}

\noindent \mpifunc{MPI\_PSEND\_INIT} 
creates a partitioned communication request and binds to it all the
arguments of a partitioned send operation.
Matching follows the same \MPI/ matching rules as for point-to-point communication (see Chapter~\ref{chap:pt2pt}) with communicator, tag, and source dictating
message matching. In the event that the communicator, tag, and source do not
uniquely identify a message, the order in which partitioned 
communication \emph{initialization} calls are made is the order in which they will eventually match.
This operation can only match with partitioned communication initialization operations,
therefore it is required to be matched with a corresponding
\mpifunc{MPI\_PRECV\_INIT} call. 
Partitioned communication initialization calls are
local. It is erroneous to provide a \mpiarg{partitions} value $\le 0$.
Send-side and receive-side buffers must be identical in size.

\begin{implementors}
  Unlike \mpifunc{MPI\_SEND\_INIT}, \mpifunc{MPI\_PSEND\_INIT} can be matched as early as the initialization call.
  Also, unlike \mpifunc{MPI\_SEND\_INIT}, \mpifunc{MPI\_PSEND\_INIT} takes an info argument.  
\end{implementors}
 
\cdeclindex{MPI\_Request}
\begin{mpi-binding}
    function_name("MPI_Precv_init")
    parameter("buf", "BUFFER", direction="IN", desc=r"initial address of recv buffer")
    parameter("partitions", "PARTITION", direction="IN", desc=r"number of partitions")
    parameter("count", "XFER_NUM_ELEM_NNI", direction="IN", desc=r"number of elements received per partition")
    parameter("datatype", "DATATYPE", direction="IN", desc=r"type of each element")
    parameter("source", "RANK", direction="IN", desc=r"rank of source")
    parameter("tag", "TAG", direction="IN", desc=r"message tag")
    parameter("comm", "COMMUNICATOR", direction="IN", desc=r"communicator")
    parameter("info", "INFO", desc=r"info argument")
    parameter("request", "REQUEST", direction="OUT", desc=r"communication request")
\end{mpi-binding}

\begin{rationale}
The \mpiarg{info} argument is provided in order to support per-operation implementation\hyphen/defined info keys.
\end{rationale}

\mpifunc{MPI\_PRECV\_INIT} creates a partitioned communication receive request
and binds to it all the
arguments of a partitioned receive operation.
This operation can only match with partitioned communication initialization operations,
therefore the \MPI/ library is required to match \mpifunc{MPI\_PRECV\_INIT} calls only with a corresponding
\mpifunc{MPI\_PSEND\_INIT} call. 
Matching follows the same \MPI/ matching rules as for point-to-point communication (see Chapter~\ref{chap:pt2pt}) with communicator, tag, and source dictating
message matching. In the event that the communicator, tag, and source do not
uniquely identify a message, the order in which partitioned 
communication initialization calls are made
is the order in which they will eventually match.
Partitioned communication initialization calls are
local. That is, \mpifunc{MPI\_PRECV\_INIT} may return before the operation
completes. It is erroneous to provide a \mpiarg{partitions} value $\le 0$. Wildcards for source and tag are not allowed.

\begin{implementors}
Unlike \mpifunc{MPI\_RECV\_INIT}, \mpifunc{MPI\_PRECV\_INIT} may communicate.
Also unlike \mpifunc{MPI\_RECV\_INIT}, \mpifunc{MPI\_PRECV\_INIT} takes an \mpiarg{info} argument.
\end{implementors}

\cdeclindex{MPI\_Request}
\begin{mpi-binding}
    function_name("MPI_Pready")
    parameter("partition", "PARTITION", direction="IN", desc=r"partition to mark ready for transfer")
    parameter("request", "REQUEST", direction="lis:inout,param:in", desc=r"partitioned communication request")
\end{mpi-binding}

\mpifunc{MPI\_PREADY} is a send-side call that indicates that a given \mpiarg{partition} is ready to be
transferred.  It is erroneous to use \mpifunc{MPI\_PREADY} on any request object that does not correspond to a partitioned send operation.
The partitioning is defined by the \mpifunc{MPI\_PSEND\_INIT} call.
Partition numbering starts at zero and ranges to one less than the number of partitions declared in
the \mpifunc{MPI\_PSEND\_INIT} call.  
Specifying a partition number that is equal to or larger than the number of partitions is erroneous.
After a call to \mpifunc{MPI\_START}/\mpifunc{MPI\_STARTALL}, all partitions associated with that operation are inactive. A call to
\mpifunc{MPI\_PREADY} marks the indicated partition as active. 
Calling \mpifunc{MPI\_PREADY} on an active partition is erroneous.

\cdeclindex{MPI\_Request}
\begin{mpi-binding}
    function_name("MPI_Pready_range")
    parameter("partition_low", "PARTITION", direction="IN", desc=r"lowest partition ready for transfer")
    parameter("partition_high", "PARTITION", direction="IN", desc=r"highest partition ready for transfer")
    parameter("request", "REQUEST", direction="lis:inout,param:in", desc=r"partitioned communication request")
\end{mpi-binding}

A call to \mpifunc{MPI\_PREADY\_RANGE} has the same effect as calls to
\mpifunc{MPI\_PREADY}, executed for \mpicode{i=partition\_low, $\ldots$,
partition\_high}, in some arbitrary order.
Calls to \mpifunc{MPI\_PREADY\_RANGE} follow the same rules as those for 
\mpifunc{MPI\_PREADY} calls.

\cdeclindex{MPI\_Request}
\begin{mpi-binding}
    function_name("MPI_Pready_list")
    parameter("length", "ARRAY_LENGTH", direction="IN", desc=r"list length")
    parameter("array_of_partitions", "PARTITION", direction="IN", desc=r"array of partitions", length="length", constant="true")
    parameter("request", "REQUEST", direction="lis:inout,param:in", desc=r"partitioned communication request")
\end{mpi-binding}

A call to \mpifunc{MPI\_PREADY\_LIST} has the same effect as calls to
\mpifunc{MPI\_PREADY}, executed for the partitions specified by the elements
\[
\mpicode{$array\_of\_partitions[0]$ ,$\ldots$, $array\_of\_partitions[count-1]$}
\]
of the \mpiarg{array\_of\_partitions}, executed in some arbitrary order.
Calls to \mpifunc{MPI\_PREADY\_LIST} follow the same rules as those for
\mpifunc{MPI\_PREADY} calls.

\subsection{Communication Completion under Partitioning}
\mpitermtitleindexmainsub{nonblocking!persistent}{partitioned completion}
\label{subsec:part-commend}

The functions \mpifunc{MPI\_WAIT} and \mpifunc{MPI\_TEST} (and
variants) are used to complete a partitioned
communication operation.  The completion of a partitioned send
operation indicates that the sender is now free to call \mpifunc{MPI\_START}/\mpifunc{MPI\_STARTALL} to restart
the operation and
subsequently \mpifunc{MPI\_PREADY}, \mpifunc{MPI\_PREADY\_RANGE} or
\mpifunc{MPI\_PREADY\_LIST}. Alternatively, the user can safely free 
the partitioned communication request after the completion of
the partitioned operation.
For the sending process, completion of the
partitioned send operation does not indicate that the
partitions of the message have all been received. 

The completion of a partitioned receive operation through \mpifunc{MPI\_WAIT} or
\mpifunc{MPI\_TEST} indicates that the receive buffer contains all of the
partitions. A function for probing the partial reception of the receive buffer is provided by 
\mpifunc{MPI\_PARRIVED}.
The \mpifunc{MPI\_PARRIVED} function can be used to determine if
the message data for the indicated partition has been received into the receive buffer. Upon success, the
receiver becomes free to access the indicated partition (as well as any others
that previously completed for that operation).  

\cdeclindex{MPI\_Request}
\begin{mpi-binding}
    function_name("MPI_Parrived")
    parameter("request", "REQUEST", direction="IN", desc=r"partitioned communication request")
    parameter("partition", "PARTITION", direction="IN", desc=r"partition to be tested")
    parameter(
        "flag",
        "LOGICAL",
        direction="out",
        desc=r"\mpicode{true} if operation completed on the specified partition, \mpicode{false} if not",
    )
\end{mpi-binding}

The function \mpifunc{MPI\_PARRIVED} can be used to
test partial completion of partitioned receive operations.
A call to \mpifunc{MPI\_PARRIVED} on an active partitioned communication request
returns \mpiarg{flag}\mpicode{ = true} if the
operation identified by \mpiarg{request} for the specified \mpiarg{partition} is complete.  
The request is not marked as complete/inactive by this procedure.
A subsequent call to an \MPI/ completing procedure (e.g., \mpifunc{MPI\_TEST}/\mpifunc{MPI\_WAIT}) is required to
complete the operation, as described in Chapter~\ref{chap:pt2pt}.
\mpifunc{MPI\_PARRIVED} may be called multiple times for a partition.
\mpifunc{MPI\_PARRIVED} may be called with a null or inactive \mpiarg{request}
argument. In either case, the operation returns with \mpiarg{flag}\code{ = true}.
Calling \mpifunc{MPI\_PARRIVED} on a request that does not correspond to a partitioned
receive operation is erroneous.

Repeated calls to \mpifunc{MPI\_PARRIVED} with the same \mpiarg{request} and \mpiarg{partition} arguments will 
eventually return \mpiarg{flag}\code{ = true} if the corresponding partitioned send operation has
been started and all send partitions have been marked as ready. For additional information on \MPI/ 
\mpitermni{progress}\mpitermindex{progress!partitioned point-to-point communication}
\mpitermindex{semantics!partitioned point-to-point communication!progress} see
\namedref{Sections}{subsec:terms:progress} \namedref{and}{subsec:pt2pt-semantics}.

\begin{implementors}
A high quality implementation will eventually return \mpiarg{flag}\mpicode{ =
true} from \mpifunc{MPI\_PARRIVED} after all of the corresponding
\mpifunc{MPI\_PREADY} calls have been made for a receive-side partition, 
even if other send partitions are not yet marked as ready. 
\end{implementors}

\subsection{Semantics of Communications in Partitioned Mode}
\mpitermtitleindex{semantics!nonblocking partitioned communications}
\label{subsec:part-semantics}

The semantics of nonblocking partitioned communication are defined by suitably extending the
definitions in Section~\ref{sec:pt2pt-semantics}.

\paraheading{Interpretation of count and datatype for partitioned communication.}
Partitioned communication uses the \mpiarg{count} and \mpiarg{datatype}
arguments in the partitioned communication initialization functions to describe
a single partition. The argument \mpiarg{partitions} specifies how many equal
partitions of a number (\mpiarg{count}) of objects of \mpiarg{datatype}s make up the
entire buffer to be transferred in the partitioned communication. As partitioned
communication describes many partitions, using absolute displacements in datatypes
(e.g., \mpiconst{MPI\_BOTTOM}) is not supported. Partitions are contiguous in memory, there
is no padding in between them.
Once a partitioned send operation is started, each partition must be marked as ready
using \mpifunc{MPI\_PREADY} and the operation must be completed using a completion
function, such as \mpifunc{MPI\_TEST} or \mpifunc{MPI\_WAIT}.


\paraheading{Order.}
Matching follows the same \MPI/ matching rules as for point-to-point communication (see Chapter~\ref{chap:pt2pt}) with communicator, tag, and source dictating
message matching. In the event that the communicator, tag, and source do not
uniquely identify the message, the order in which partitioned 
communication initialization calls are made is the order in which they will eventually match.

\section{Partitioned Communication Examples}
\label{sec:part-communication-examples}
This section provides concrete examples of the utility of partitioned communication
in realistic settings.
\subsection{Partition Communication with Threads/Tasks Using OpenMP 4.0 or later}
\label{subsec:part-channel-examples}
In the following the equal partitioning on send-side and receive-side Example~\ref{part-example1} is extended
to utilize threads. In this case, the receive-side uses the same
number of partitions as the send-side as in the previous example, but this
example uses multiple threads on the send-side. Note that the
\mpifunc{MPI\_PSEND\_INIT} and \mpifunc{MPI\_PRECV\_INIT} functions match
each other like in the previous example.

\begin{example}
Equal partitioning on send-side and receive-side using threads.
\label{part-example2}%
\exindex{Partitioned communication!Equal send/recv partitioning}%
\Exindex{C}{Partitioned communication using threads}{MPI\_Psend\_init,MPI\_Precv\_init,MPI\_Pready}%
%%HEADER
%%LANG: C
%%DECL: void error(const char *);
%%ENDHEADER
\begin{lstlisting}[language={[MPI]C}]
#include <stdlib.h>
#include "mpi.h"
#define NUM_THREADS 8
#define PARTITIONS 8
#define PARTLENGTH 16
int main(int argc, char *argv[]) /* same send/recv partitioning */
{
  double message[PARTITIONS*PARTLENGTH];
  int partitions = PARTITIONS;
  int partlength = PARTLENGTH;
  int count = 1, source = 0, dest = 1, tag = 1, flag = 0;
  int myrank;
  int provided;
  MPI_Request request;
  MPI_Info info = MPI_INFO_NULL;
  MPI_Datatype xfer_type;
  MPI_Init_thread(&argc, &argv, MPI_THREAD_MULTIPLE, &provided);
  if (provided < MPI_THREAD_MULTIPLE)
     MPI_Abort(MPI_COMM_WORLD, EXIT_FAILURE);
  MPI_Comm_rank(MPI_COMM_WORLD, &myrank);
  MPI_Type_contiguous(partlength, MPI_DOUBLE, &xfer_type);
  MPI_Type_commit(&xfer_type);
  if (myrank == 0)
  {
     MPI_Psend_init(message, partitions, count, xfer_type, dest, tag,
                    MPI_COMM_WORLD, info, &request);
     MPI_Start(&request);

     #pragma omp parallel for shared(request) num_threads(NUM_THREADS)
     for (int i=0; i<partitions; i++)
     {
        /* compute and fill partition #i, then mark ready: */
        MPI_Pready(i, request);
     }
     while(!flag)
     {
        /* Do useful work */
        MPI_Test(&request, &flag, MPI_STATUS_IGNORE);
        /* Do useful work */
     }
     MPI_Request_free(&request);
  }
  else if (myrank == 1)
  {
     MPI_Precv_init(message, partitions, count, xfer_type, source, tag,
                    MPI_COMM_WORLD, info, &request);
     MPI_Start(&request);
     while(!flag)
     {
        /* Do useful work */
        MPI_Test(&request, &flag, MPI_STATUS_IGNORE);
        /* Do useful work */
     }
     MPI_Request_free(&request);
  }
  MPI_Finalize();
  return 0;
}
\end{lstlisting}
\end{example}

\subsection{Send-only Partitioning Example with Tasks and OpenMP version 4.0 or later}
The previous example is tailored specifically for send-side partitioning using threads.
This is an example where parallel task producers produce
input to part of an overall buffer; they complete in any order and contribute
to the overall buffer. 
\begin{example}
Parallel task producers for partitioned communication using threads.
\label{partexample3}%
\exindex{Partitioned communication!Send with tasks}%
\Exindex{C}{Partitioned communication with tasks}{MPI\_Psend\_init,MPI\_Precv\_init,MPI\_Pready}%
%%HEADER
%%LANG: C
%%DECL: void error(const char *);
%%ENDHEADER
\begin{lstlisting}[language={[MPI]C},basicstyle=\lstsmall]
#include <stdlib.h>
#include "mpi.h"
#define NUM_THREADS 8
#define NUM_TASKS 64
#define PARTITIONS NUM_TASKS
#define PARTLENGTH 16
#define MESSAGE_LENGTH PARTITIONS*PARTLENGTH
int main(int argc, char *argv[]) /* send-side partitioning */
{
  double message[MESSAGE_LENGTH];
  int send_partitions = PARTITIONS,
      send_partlength = PARTLENGTH,
      recv_partitions = 1,
      recv_partlength = PARTITIONS*PARTLENGTH;
  int count = 1, source = 0, dest = 1, tag = 1, flag = 0;
  int myrank;
  int provided;
  MPI_Request request;
  MPI_Info info = MPI_INFO_NULL;
  MPI_Datatype send_type;
  MPI_Init_thread(&argc, &argv, MPI_THREAD_MULTIPLE, &provided);
  if (provided < MPI_THREAD_MULTIPLE)
     MPI_Abort(MPI_COMM_WORLD, EXIT_FAILURE);
  MPI_Comm_rank(MPI_COMM_WORLD, &myrank);
  MPI_Type_contiguous(send_partlength, MPI_DOUBLE, &send_type);
  MPI_Type_commit(&send_type);

  if (myrank == 0)
  {
     MPI_Psend_init(message, send_partitions, count, send_type, dest, tag,
                    MPI_COMM_WORLD, info, &request);
     MPI_Start(&request);

     #pragma omp parallel shared(request) num_threads(NUM_THREADS)
     {
        #pragma omp single
        {
           /* single thread creates 64 tasks to be executed by 8 threads */
           for (int partition_num=0;partition_num<NUM_TASKS;partition_num++)
           {
              #pragma omp task firstprivate(partition_num)
              {
                 /* compute and fill partition #partition_num, then mark
                 ready: */
                 /* buffer is filled in arbitrary order from each task */
                 MPI_Pready(partition_num, request);
              } /*end task*/
           } /* end for */
        } /* end single */
     } /* end parallel */
     while(!flag)
     {
        /* Do useful work */
        MPI_Test(&request, &flag, MPI_STATUS_IGNORE);
        /* Do useful work */
     }
     MPI_Request_free(&request);
  }
  else if (myrank == 1)
  {
     MPI_Precv_init(message, recv_partitions, recv_partlength, MPI_DOUBLE,
                    source, tag, MPI_COMM_WORLD, info, &request);

     MPI_Start(&request);
     while(!flag)
     {
        /* Do useful work */
        MPI_Test(&request, &flag, MPI_STATUS_IGNORE);
        /* Do useful work */
     }
     MPI_Request_free(&request);
  }
  MPI_Finalize();
  return 0;
}
\end{lstlisting}
\end{example}

\subsection{Send and Receive Partitioning Example with OpenMP version 4.0 or later}
This example demonstrates receive-side partial completion notification
using more than one partition per receive-side thread. It uses a naive flag 
based method to test for multiple completed partitions per thread.
Note that this means that some threads may be busy polling for completion of
assigned partitions when partitions are available to work on that were not 
assigned to the polling threads in this example. More advanced work stealing methods
could be employed for greater efficiency.
Like previous examples, it also demonstrates send-side production of
input to part of an overall buffer. This example also uses different send-side and
receive-side partitioning. 

\begin{example}
\label{partexample4}
\exindex{Partitioned communication!Partial completion notification}%
\Exindex{C}{Partitioned communication with partial completion}{MPI\_Parrived,MPI\_Psend\_init,MPI\_Precv\_init,MPI\_Pready}%
Partitioned communication receive-side partial completion.\\
%%HEADER
%%LANG: C
%%DECL: void error(const char *);
%%ENDHEADER
\begin{lstlisting}[language={[MPI]C}]
#include <stdlib.h>
#include "mpi.h"
#define NUM_THREADS 64
#define PARTITIONS NUM_THREADS
#define PARTLENGTH 16
#define MESSAGE_LENGTH PARTITIONS*PARTLENGTH
int main(int argc, char *argv[]) /* send-side partitioning */
{
  double message[MESSAGE_LENGTH];
  int send_partitions = PARTITIONS,
      send_partlength = PARTLENGTH,
      recv_partitions = PARTITIONS*2,
      recv_partlength = PARTLENGTH/2;
  int source = 0, dest = 1, tag = 1, flag = 0;
  int myrank;
  int provided;
  MPI_Request request;
  MPI_Info info = MPI_INFO_NULL;
  MPI_Datatype send_type;
  MPI_Init_thread(&argc, &argv, MPI_THREAD_MULTIPLE, &provided);
  if (provided < MPI_THREAD_MULTIPLE)
     MPI_Abort(MPI_COMM_WORLD, EXIT_FAILURE);
  MPI_Comm_rank(MPI_COMM_WORLD, &myrank);
  MPI_Type_contiguous(send_partlength, MPI_DOUBLE, &send_type);
  MPI_Type_commit(&send_type);

  if (myrank == 0)
  {
     MPI_Psend_init(message, send_partitions, 1, send_type, dest, tag,
                    MPI_COMM_WORLD, info, &request);
     MPI_Start(&request);
     #pragma omp parallel for shared(request) \
                              firstprivate(send_partitions) \
                              num_threads(NUM_THREADS)
     for (int i=0; i<send_partitions; i++)
     {
        /* compute and fill partition #i, then mark ready: */
        MPI_Pready(i, request);
     }
     while(!flag)
     {
        /* Do useful work */
        MPI_Test(&request, &flag, MPI_STATUS_IGNORE);
        /* Do useful work */
     }
     MPI_Request_free(&request);
  }
  else if (myrank == 1)
  {
     MPI_Precv_init(message, recv_partitions, recv_partlength,
                    MPI_DOUBLE, source, tag, MPI_COMM_WORLD, info,
                    &request);
     MPI_Start(&request);
     #pragma omp parallel for shared(request) \
                              firstprivate(recv_partitions) \
                              num_threads(NUM_THREADS)
     for (int j=0; j<recv_partitions; j+=2)
     {
        int part_flag = 0;
        int part1_complete = 0;
        int part2_complete = 0;
        while(part1_complete == 0 || part2_complete == 0)
        {
           /* test partition #j and #j+1 */
           MPI_Parrived(request, j, &part_flag);
           if(part_flag && part1_complete == 0)
           {
              part1_complete++;
              /* Do work using partition j data */
           }
           MPI_Parrived(request, j+1, &part_flag);
           if(part_flag && part2_complete == 0)
           {
              part2_complete++;
              /* Do work using partition j+1 */
           }
        }
     }
     while(!flag)
     {
        /* Do useful work */
        MPI_Test(&request, &flag, MPI_STATUS_IGNORE);
        /* Do useful work */
     }
     MPI_Request_free(&request);
  }
  MPI_Finalize();
  return 0;
}
\end{lstlisting}
\end{example}
